\section{从 checkpoint 到 solution：最后一公里才是系统主体}

课堂反复强调，绝大多数企业拿到 pretrained checkpoint 后仍无法直接使用。权重缺少稳定 runtime、服务接口、认证、数据连接、task-specific evaluation、更新策略和故障处理。Mistral 的商业定位因此不是只“卖一个模型”，而是帮助客户把模型嵌入 on-prem、private cloud 或 edge 工作流。

\subsection{五层产品栈}

先看从权重到运营的五层结构。每层都可能单独成为失败点：runtime OOM、endpoint 限流、tool 调用错误、应用缺少权限边界、运营没有回滚。读这张图时要从右向左追问：如果最终业务结果失败，系统是否能沿着 operations、application、endpoint 和 runtime 定位到具体责任层，而不是把所有问题都归因于“模型不够聪明”。

\lecturefigure{09-checkpoint-to-solution.png}{从 model weights 到 production operations 的五层差距。}{本地字幕 20:25--24:30；概念重绘。}

\begin{knowledgebox}{读图：为什么稳定 HTTP endpoint 已经是产品能力}
一个 endpoint 必须处理并发、超时、排队、版本、认证、配额、日志和错误语义。即使模型本身不变，batching policy 或 cache 策略也会改变延迟分布。客户看到的是完整服务的可用性，而不是离线 benchmark。
\end{knowledgebox}

\subsection{部署模式由约束选择}

\term{on-prem} 指模型运行在客户自有设施；\term{private cloud} 指在隔离的云环境中运行；\term{edge} 指接近设备或现场、资源和连接受限的部署。三者都不是“比 API 更高级”，而是在数据驻留、控制、延迟、升级速度和成本之间选择不同位置。下面比较的重点不是产品名称，而是每种模式把哪些责任留给 provider、哪些责任转移给 customer，以及这种责任分配是否匹配组织能力。

\lecturefigure{10-deployment-modes.png}{API、private cloud、on-prem 和 edge 是不同约束组合。}{本地字幕 20:25--24:30、32:00--36:45；概念重绘。}

\begin{warningbox}{私有部署不自动更安全}
把模型放进客户网络可以增强数据控制，却同时把 patching、访问控制、密钥、监控和 incident response 责任交给客户或供应商。安全来自明确责任和可验证控制，不来自“服务器在自己机房”这一位置事实。
\end{warningbox}

\subsection{定制不是一次 fine-tune}

课堂给出的实际流程包括理解 use case、收集少量真实样本、生成 synthetic data、fine-tune、评测并部署。少量高质量样本可以显著改善行为，但前提是系统知道“什么算改善”，并防止通用能力、安全性或格式遵循发生回退。下面的闭环把定制视为持续控制过程：每一次部署都会产生新失败证据，而这些证据必须先进入评测，再决定是否生成数据和更新权重。

\lecturefigure{11-customization-loop.png}{企业定制必须由评测闭环驱动，而不是训练完一次就结束。}{本地字幕 22:40--24:30；概念重绘。}

\begin{importantbox}{定制闭环的最小交付物}
\begin{enumerate}
  \item 任务定义与失败 taxonomy；
  \item 可追溯的真实/合成训练数据；
  \item task metric、regression suite 与安全测试；
  \item 可复现的训练配置和 base model 版本；
  \item serving 监控、回滚条件与数据回流规则。
\end{enumerate}
没有这些交付物，fine-tuning 只是不可审计的权重变化。
\end{importantbox}

\teachervoice{老师提醒，即使最技术化的企业也希望供应商工程师与业务团队深度协作，因为真正困难的是把领域判断转成数据、指标和系统边界。这个观察解释了为什么“模型能力商品化”并不会消灭解决方案工程，反而会把价值推向集成层。}

\subsection{本章小结}

Checkpoint-to-solution 需要 runtime、接口、应用、运营和定制闭环。部署模式没有统一优劣；应先列出不可妥协的隐私、可靠性、延迟和控制要求，再选择架构。

